\documentclass{article}
\usepackage{graphicx} % Required for inserting images
\usepackage{float}

\usepackage{float} 

\usepackage{natbib}

\usepackage{amssymb}
\usepackage{amsmath}
\usepackage{amsthm}
\usepackage{enumitem}

\usepackage{tcolorbox}
\usepackage{enumitem}
\tcbuselibrary{breakable, skins}

\usepackage{xcolor}
\usepackage{graphicx}
\usepackage{wrapfig}
\usepackage[makeroom]{cancel}
\newcommand{\red}[1]{\textcolor{red}{#1}}
\newtheorem{proposition}{Proposition}
\usepackage{bm} 

\newtheorem{definition}{Definition}



\title{Relative Multiplicity}
\author{minigonche }
\date{September 2026}

\begin{document}

\maketitle

\section{Introduction}

Although useful in many contexts, distance-based multiplicity ($\delta M_\sigma$) is inadequate for a particular but recurrent case in ecology. By its construction, one can define the value of $\sigma$ which will numerically encode the notion of: ``two totally dissimilar units or subunits". This is fine in specific cases where all biological entities are numerically comparable, but what if we need to compare two actually \textbf{completely} different entities? Like for example comparing CAZymes, peptidases and lipases? Or even enzymes form the same group but different families? Ecologist agree that these are different entities, but effectively selecting a value for $\sigma$ means that ecologist will need to agree on \textit{how} dissimilar, numerically speaking, these enzymatic groups actually are. \\

Another limitation that $\delta M_\sigma$ inherit from traditional biodiversity indexes, is that ecological relevance is intrinsically encoded as abundance. Again, useful in specific cases where the biological units are comparable. But if one is studying together ants and felines, equating the ecological relevance of each of these entities with their shier numbers seems unfair, to say the least.\\

These two limitations motivate the construction of an index that addresses the following scenario:


\begin{quote}
Given a sample of truly incomparable biological units, how diverse are each of the units, with respect to themselves, in average?
\end{quote}

Clearly this index would be the highest when the diversity inside each unit is maximum, but that maximum depends entirely on the particulars of each unit. Here is a concrete simple example. Suppose we have a sample with three units: PL10, AA10 and GT51 all of them CAZYme families from different catalytic groups. According to the database, their total number of different sequences are roughly: 1.5k, 13k and 110k, respectively. This is how the index should respond to different scenarios:

\begin{table}[H]
\centering
\begin{tabular}{l|cc|cc|cc|c}
\hline
 & \multicolumn{2}{c|}{PL10 (1.5k)} & \multicolumn{2}{c|}{AA10 (13k)} & \multicolumn{2}{c|}{GT51 (110k)} & \\
Scenario & Distinct & Score & Distinct & Score & Distinct & Score & Index \\
\hline
A & 750     & 0.50 & 6{,}500 & 0.50 & 55{,}000  & 0.50 & 0.50 \\
B & 1{,}500 & 1.00 & 1{,}500 & 0.12 & 1{,}500   & 0.01 & 0.38 \\
C & 0       & 0.00 & 0       & 0.00 & 110{,}000 & 1.00 & 0.33 \\
\hline
\end{tabular}
\label{tab:scenarios}
\end{table}

Here the index is simply averaging the individual scores. Although trivial, it shows the intended motivation: each internal diversity score is defined by the unit's theoretical maximum diversity (some have naturally higher diversity than others) and the index itself is not influenced by the fact that GT51 might probably have 70 times more subunits (and maybe even abundance) than PL10.

This is an example in the realm of inventory diversity and Richness, but we can extend this idea to abundance aware and distance-based diversity. 

\section{Relative Multiplicity}

This is the motivation behind \textit{Relative Multiplicity}. A normalized, distance-based diversity measure that tracks how diverse are the the units in an ensemble internally, relative to themselves.

Let us start with a set of biological units, for example proteins, genes or genera. For relative multiplicity to have meaning, we need each of the units to have two thing:

\begin{itemize}
    \item Internal members, which we will refer to as subunits. This should be biological entities that the unit is grouping together. Following our example, this could be variants for proteins, alleles for genes or species for genera. This is what enables us to talk about how \textit{diverse} a unit is, which is what multiplicity measures.
    \item A $referenence$ composition. Each unit should have a list of subunits together, with their abundance distribution, that acts as the baseline composition against which the actual state of the unit is compared. This could be all registered variants of a given protein in a database, all known alleles in a specific population or the geographic relevant species for a genera. The sample's diversity will be \textit{relative} to this reference diversity.
\end{itemize}


For its construction, we will use $\delta D_\sigma$ as defined in \citep{GonzalezCasabianca2026multiplicity}, but the construction is analogous to any other alpha diversity measure. Recall its definition for a unit of $S$ subunits, $P$ its relative abundance vector and $\Delta$ its associated distance matrix. Then the functional diversity for the unit is: 

$$
\delta D_\sigma =  \left( 1-\sum\limits_{i,j = 1}^S \frac{min(\sigma, \Delta_{ij})}{\sigma} P_iP_j \right)^{-1} = \left(1-\frac{Q_\sigma}{\sigma} \right)^{-1}
$$

where $Q_\sigma = P^t \Delta_\sigma P$ and $(\Delta_\sigma)_{ij} = min(\sigma, \Delta_{ij})$. Now for relative multiplicity:

\begin{definition}
  Given a sample of $k$ units, its relative multiplicity is:

$$
RM = \frac{1}{\left|\{ f : N_f > 0 \}\right|} \sum_{f \,:\, N_f > 0} \frac{\delta D_\sigma (f)}{\delta D_\sigma (\tilde{f})}
$$
  
  \label{def:relative_multipplicity}
\end{definition}

were $\tilde{f}$ is the reference composition of the unit, $f$ its  actual composition and $N_f$ its actual abundance.






\bibliographystyle{plainnat}
\bibliography{biblio}











\end{document}
